\chapter{Introduction}

Steininger and Yurkevich~\cite{polyhedron.without.rupert} constructed a convex
polyhedron, the \emph{Noperthedron}, and proved that it does not have Rupert's
property: no copy of it can be passed through a hole in another copy. This blueprint
describes a Lean proof of the same theorem that is checked by the Lean kernel alone,
using only the standard axioms.

The overall strategy agrees with~\cite{polyhedron.without.rupert}: reduce to purely
rotational passages, use the symmetries of the Noperthedron to shrink the space of
rotation pairs to a bounded box, cover that box by a large tree of sub-boxes, and
exclude each sub-box by exhibiting an inner vertex whose shadow lies outside the outer
shadow. The details differ in several ways.

\begin{itemize}
\item \textbf{Rational coordinates.} Rotations are parametrized by Cayley
  coordinates (Chapter~\ref{chap:cayley}) instead of angles. Every rotation matrix
  in the case analysis is a rational function of the box coordinates, so every leaf
  check is exact polynomial arithmetic. The only irrational numbers are the fixed
  constants $\cos(2\pi k/15)$, $\sin(2\pi k/15)$ and the vertices
  (Chapter~\ref{chap:data}).
\item \textbf{A different reduction.} The symmetry reduction
  (Chapter~\ref{chap:reduction}) puts the outer view in a wedge, picks the relative
  rotation from its orbit under the symmetries by maximizing a score, and covers the
  result by five Cayley charts.
\item \textbf{No local theorem.} Near the configurations where the inner and outer
  shadows coincide, \cite{polyhedron.without.rupert} uses a separate local argument.
  Here these configurations are blown up (Section~\ref{sec:blowup}): the vanishing
  gap factors exactly, and the remaining factor is checked like any other polynomial.
\item \textbf{Polynomial bounds instead of Lipschitz bounds.} A box is excluded by
  a lower bound of a polynomial on $[-1,1]^5$ (Chapter~\ref{chap:polybounds}), computed
  from its coefficients or from its Bernstein coefficients.
\end{itemize}

The case analysis is a tree with $245{,}484$ leaves. It is produced by an untrusted
floating-point search and checked by a Lean function that the kernel evaluates
(Chapter~\ref{chap:main}).
